\section{Proposed Fuzzy PI Controller}
\label{sec:fpi_controller}

\subsection{Local measurements and continuous-clock interface}
\label{subsec:fpi_laws}
Let $\mathcal N_i^-$ be the fixed set of incoming neighbors and
$d_i=|\mathcal N_i^-|>0$. For common capacity $C$ and target
$\beta_{ji}^*$, define frame-sum and percentage errors
\begin{equation}
 E_{i,k}=\sum_{j\in\mathcal N_i^-}(\beta_{ji,k}-\beta_{ji}^*),
 \quad e_{i,k}=-\gamma_iE_{i,k},\quad\gamma_i=\frac{100}{d_iC}.
 \label{eq:fpi_error}
\end{equation}
Thus positive $E$ calls for faster local consumption. The index $k$
is local to each node; $e$ is target minus measured average occupancy
in percentage points. Unequal capacities require consistent weighted
errors, rather than the scalar calibration below.

With local polling period $p$ and execution delay $d_{\rm ctrl}$,
measurement and application times satisfy
\begin{align}
 \theta_i(t^m_{i,k})&=\theta_i(0)+kp,\nonumber\\
 \theta_i(t^a_{i,k})&=\theta_i(0)+kp+d_{\rm ctrl}.
 \label{eq:fpi_local_schedule}
\end{align}
The product $kp$ is update index times period, not gain $k_p$.
Requested correction $q_{i,k}$ is limited to
$u_{i,k}=\operatorname{sat}_{u_{\max}}q_{i,k}$ and held until the
next application, with $\dot\theta_i=\omega_i^u+u_{i,k}$.
The physical measurement-to-actuation age therefore depends on the
past oscillator rate. Initial correction is zero.

Actual arrivals use compatible negative-time phase history and an
order-preserving delivery map. For desired send-time latency
$\tau_{ji}(s)$, arrival time is
$a_{ji}(s)=\sup_{r\le s}\{r+\tau_{ji}(r)\}$; counts follow delivered
integer send crossings, not an arbitrary continuous cutoff. This
permits ordered batch release and imposes no finite serialization rate.
Appendix~\ref{app:controller} specifies count initialization, the
in-flight offset and the complete execution interface.

\subsection{Fuzzification, inference and control update}
The implemented integrate-first PI baseline uses sample sum $I$, fixed integration
scale $h=p/\omega_{\rm base}$, and $\kappa_I=k_ih$:
\begin{equation}
 I_k=I_{k-1}+E_k,\qquad q_k^{\rm PI}=k_pE_k+\kappa_I I_k.
 \label{eq:fpi_native}
\end{equation}
The sum has no elapsed-time multiplier. The proposed Fuzzy law is
\begin{align}
 X_k&=K_e e_k,\quad Y_k=K_{\Delta e}(e_k-e_{k-1}),\nonumber\\
 A_k&=A_{k-1}+K_uF(X_k,Y_k),\qquad q_k^{\rm F}=A_k,
 \label{eq:fpi_fuzzy_law}
\end{align}
where $K_e,K_{\Delta e}>0$, $K_u<0$ and $A$ stores the requested
correction. The second input is a sample difference, not a time
derivative. All rules and gains remain fixed during a run.

Fuzzification uses five labels: negative large (NL), negative small (NS), zero (ZE),
positive small (PS) and positive large (PL). Their input centers are
$(-10,-5,0,5,10)$. Interior memberships are triangular:
$\varphi_a(z)=[1-|z-c_a|/5]_+$, $a=2,3,4$, where
$[z]_+=\max(z,0)$. The outer shoulders are
\begin{equation}
 \varphi_1(z)=\operatorname{clip}_{[0,1]}\frac{-z-5}{5},\quad
 \varphi_5(z)=\operatorname{clip}_{[0,1]}\frac{z-5}{5}.
\end{equation}
Here $\operatorname{clip}_{[0,1]}$ projects onto $[0,1]$.
Consequent singleton values are $s=(-5,-2.5,0,2.5,5)$.
Inference uses Mamdani minimum rule activation and singleton
weighted-average defuzzification. All 25 rules are specified by
$\ell(a,b)=\min(5,\max(1,a+b-3))$ and
\begin{equation}
 w_{ab}=\min(\varphi_a(X),\varphi_b(Y)),\quad
 F=\frac{\sum_{a,b=1}^5w_{ab}s_{\ell(a,b)}}{\sum_{a,b=1}^5w_{ab}}.
 \label{eq:fpi_rules}
\end{equation}
Rules sharing a consequent are summed separately; replacing their
sum by maximum aggregation changes the controller. The denominator
is positive for finite inputs, and $|F|\le5$ bounds each increment,
not the indefinitely accumulated correction. Initialize
$I_{-1}=A_{-1}=e_{-1}=0$. External clipping retains raw $A$;
anti-windup, observer and mode switching are disabled.
Table~\ref{tab:fpi_full_rules} also gives the complete consequent matrix.
This specification fixes the implemented nonlinear map; fuzzy inference
is not used here as an unspecified online gain tuner.

\begin{table}[t]
\caption{All consequent labels of the proposed Fuzzy law. Rows and columns
are the present-error and error-difference inputs, respectively.}
\label{tab:fpi_full_rules}
\centering
\begin{tabular}{@{}c|ccccc@{}}\toprule
$X\backslash Y$ & NL & NS & ZE & PS & PL\\\midrule
NL & NL & NL & NL & NS & ZE\\
NS & NL & NL & NS & ZE & PS\\
ZE & NL & NS & ZE & PS & PL\\
PS & NS & ZE & PS & PL & PL\\
PL & ZE & PS & PL & PL & PL\\\bottomrule
\end{tabular}
\end{table}

\subsection{PI calibration and nonlinear mechanism}
\label{subsec:fpi_calibration}
For $F_0(X,Y)=(X+Y)/2$, coefficient matching with
$\Delta q_k^{\rm PI}=k_p\Delta E_k+\kappa_I E_k$ gives
\begin{equation}
 k_p=-\frac{K_u\gamma_iK_{\Delta e}}2,\qquad
 \kappa_I=-\frac{K_u\gamma_iK_e}2.
 \label{eq:fpi_matching}
\end{equation}
\begin{btdproposition}[Calibrated linear counterpart]
\label{prop:fpi_equivalence}
With~\eqref{eq:fpi_matching}, common measurements, consistent zero
constructors and identical execution policies, replacing $F$ by $F_0$
gives the PI baseline's requested correction at every update, including
$q_0=(k_p+\kappa_I)E_0$.
\end{btdproposition}
This is a scale and implementation equivalence for $F_0$, not full
nonlinear equivalence. The substitution and initialization proof are
in Appendix~\ref{app:controller}.

This mapping controls a common comparison pitfall: matching gain names
without matching the measured error, integration order and time scale
does not match controllers. For the triangle, $\gamma=100/4096$ and
$h=1000/(29999/30000)$, so the moderate coordinates
$(k_p,k_i)=(2\times10^{-6},5\times10^{-11})$ correspond to
$K_e=0.1$, $K_{\Delta e}\simeq3.9998667$ and
$K_u\simeq-4.0961365\times10^{-5}$. Node indices were suppressed
in the update equations. On unequal-degree graphs, node-dependent
$K_{u,i}$ can compensate $\gamma_i$ when common calibrated coordinates
are required. Target and capacity conventions must stay consistent.

\begin{btdproposition}[Central nonlinear map]
\label{prop:fpi_central}
On $\Omega_r=\{(X,Y):|X|+|Y|\le r\}$, $0<r\le5$, set
$m=\min(|X|,|Y|)$ and
\begin{equation}
 J=\begin{cases}\operatorname{sgn}(X)m,&XY>0,\\0,&XY\le0.\end{cases}
 \qquad F(X,Y)=\frac{F_0(X,Y)+J}{1+2m/5}.
 \label{eq:fpi_central}
\end{equation}
There exists $\alpha(X,Y)$ with $F=\alpha F_0$ and
$\underline\alpha=(1+r/5)^{-1}\le\alpha\le2=\overline\alpha$.
\end{btdproposition}
The four active min-rule weights yield this identity. Same-sign
inputs strengthen $F_0$; opposing signs attenuate its magnitude.
Uniformly increasing PI gains cannot produce both effects. This
axis calibration is not an ordinary Jacobian linearization; details
and sector bounds are in Appendix~\ref{app:controller}.

With $\mathcal R_F=F-F_0$, the zero-constructor raw correction obeys
\begin{equation}
 A_k=k_pE_k+\kappa_I\sum_{\ell=0}^kE_\ell
               +K_u\sum_{\ell=0}^k\mathcal R_F(X_\ell,Y_\ell).
 \label{eq:fpd_accounting}
\end{equation}
The last term exposes nonlinear memory. It is accounting on Fuzzy's
own measurements, not a causal decomposition of two independently
evolving closed loops. Moreover, same-sign inputs do not always mean
growing absolute error. Performance requires closed-loop evidence.
